% =====================================================================
%  Bien dich: pdfLaTeX (mac dinh cua Overleaf) -- khong can chinh gi.
%  Neu muon dung font he thong thi doi compiler sang XeLaTeX va thay
%  ba dong inputenc/fontenc/lmodern ben duoi bang:
%      \usepackage{fontspec}
%      \setmainfont{TeX Gyre Termes}
% =====================================================================
\documentclass[11pt,a4paper]{article}

\usepackage[utf8]{inputenc}
\usepackage[T5]{fontenc}
\usepackage{lmodern}

\usepackage[margin=2.5cm]{geometry}
\usepackage{amsmath}
\usepackage{amssymb}
\usepackage{booktabs}
\usepackage{array}
\usepackage{caption}
\usepackage{microtype}
\usepackage{enumitem}
\usepackage{framed}
\usepackage{tikz}
\usetikzlibrary{arrows.meta, positioning}
\usepackage[hidelinks]{hyperref}

\renewcommand{\tablename}{Bảng}
\renewcommand{\figurename}{Hình}
\renewcommand{\contentsname}{Nội dung}
\renewcommand{\refname}{Tài liệu tham khảo}

\captionsetup[table]{skip=6pt}
\captionsetup[figure]{skip=8pt}
\setlist[itemize]{itemsep=3pt, topsep=4pt}
\setlist[enumerate]{itemsep=4pt, topsep=4pt}

\newcolumntype{L}{>{\raggedright\arraybackslash}p{6.6cm}}
\newenvironment{luuy}{\begin{leftbar}\noindent\ignorespaces}{\end{leftbar}}

\title{\bfseries Đọc hiểu RanPAC\\[2pt]
\large\mdseries Chiếu ngẫu nhiên và mô hình tiền huấn luyện cho học liên tục}
\author{Tóm tắt và diễn giải bài báo của McDonnell và cộng sự, NeurIPS 2023}
\date{Tháng 8, 2026}

\begin{document}
\maketitle

\noindent
Tài liệu này diễn giải bài báo RanPAC \cite{ranpac} cho người mới bắt đầu: bài
giải quyết vấn đề gì, ý tưởng cốt lõi là gì, từng thành phần hoạt động ra sao, và
vì sao nó liên quan tới đề tài của chúng ta. Mọi số liệu trích dẫn đều lấy trực
tiếp từ bài báo, có ghi rõ nguồn là bảng nào.

\tableofcontents
\bigskip

% =====================================================================
\section{Bài báo giải quyết vấn đề gì}

Bối cảnh là \textbf{học liên tục} (continual learning): dữ liệu đến theo từng
đợt, mô hình phải học đợt mới mà không được xem lại đợt cũ, và không được lưu ảnh
cũ. Trở ngại trung tâm là \textbf{quên thảm khốc} -- cập nhật trọng số cho dữ
liệu mới sẽ ghi đè lên những gì đã học.

Từ khi có các mô hình nền tảng lớn được tiền huấn luyện sẵn, lĩnh vực này tách
thành ba hướng.

\begin{table}[htbp]
\centering
\caption{Ba hướng khai thác mô hình tiền huấn luyện cho học liên tục.}
\label{tab:huong}
\begin{tabular}{lLL}
\toprule
Hướng & Cách làm & Điểm yếu \\
\midrule
Prompting
  & Học một túi prompt nhỏ gắn vào đầu vào của transformer (L2P, DualPrompt,
    CODA-Prompt)
  & Chỉ áp dụng được cho transformer \\
\addlinespace
Tinh chỉnh backbone
  & Tinh chỉnh nhẹ toàn mạng với tốc độ học thấp (SLCA)
  & Vẫn quên tích lũy, tốn tài nguyên huấn luyện \\
\addlinespace
Class-Prototype (CP)
  & Lấy trung bình đặc trưng của mỗi lớp làm ``đại diện'', phân loại bằng độ
    tương đồng cosine (ADaM)
  & Bị coi là đơn giản nhưng kém chính xác \\
\bottomrule
\end{tabular}
\end{table}

RanPAC thuộc hướng thứ ba. Luận điểm của bài là: \textbf{hướng CP bị đánh giá
thấp oan -- không phải vì bản chất nó yếu, mà vì người ta đang dùng nó sai cách.}

% =====================================================================
\section{Chẩn đoán: vì sao class-prototype đang kém}

Đây là phần thuyết phục nhất của bài, vì họ không suy đoán mà đo trực tiếp.

Lấy ViT-B/16 tiền huấn luyện trên ImageNet-21K, chạy trên ImageNet-R, rồi tính
\textbf{hệ số tương quan Pearson giữa từng cặp class-prototype}. Kết quả: các
prototype của những lớp \emph{khác nhau} lại tương quan với nhau rất mạnh, cao
hơn hai lần so với khi làm y hệt trên chính ImageNet-1K.

Diễn giải: khi đem mô hình sang một miền dữ liệu mới, \textbf{mọi đặc trưng đều
bị dồn về cùng một hướng}. Prototype của hai lớp hoàn toàn khác nhau trở nên
giống nhau một cách bất thường.

Hậu quả đo được ngay trên histogram độ tương đồng: phân bố ``giống lớp đúng của
mình'' và ``giống 199 lớp sai'' chồng lấn nặng.

\begin{table}[htbp]
\centering
\caption{Độ chính xác trên tập huấn luyện ImageNet-R, theo Hình 2 của bài báo.
Con số cho thấy vấn đề nằm ở cách chấm điểm, không ở bản thân ý tưởng prototype.}
\label{tab:chandoan}
\begin{tabular}{lr}
\toprule
Cách chấm điểm & Độ chính xác huấn luyện \\
\midrule
Class-prototype + cosine (NCM)              & 64\,\% \\
Class-prototype + nghịch đảo Gram, công thức \eqref{eq:gram}
                                            & 75\,\% \\
Linear probe huấn luyện chung trên cả 200 lớp & 82\,\% \\
Công thức \eqref{eq:gram} kèm chiếu ngẫu nhiên, $M = 2000$
                                            & \textbf{87\,\%} \\
\bottomrule
\end{tabular}
\end{table}

Hàng cuối đáng chú ý: nó \emph{vượt} cả linear probe được huấn luyện chung trên
toàn bộ dữ liệu. Hai mục tiếp theo giải thích hai bước đưa 64\,\% lên 87\,\%.

% =====================================================================
\section{Ý tưởng thứ nhất: khử tương quan bằng ma trận Gram}

\subsection{Đổi công thức chấm điểm}

Thay vì cosine giữa đặc trưng và prototype trung bình $\bar c_y$:
\begin{equation}
s_y = \frac{f_{\text{test}}^{\top} \bar c_y}
           {\|f_{\text{test}}\| \cdot \|\bar c_y\|},
\label{eq:cosine}
\end{equation}
bài dùng nghịch đảo ma trận Gram $G$ của đặc trưng, với $c_y$ là prototype
\emph{không} chia trung bình:
\begin{equation}
s_y = f_{\text{test}}^{\top} G^{-1} c_y .
\label{eq:gram}
\end{equation}

\subsection{Trực giác: tẩy trắng rồi mới so sánh}

Ta đi tìm một phép biến đổi tuyến tính $D$ sao cho sau khi biến đổi, ma trận Gram
trở thành ma trận đơn vị -- nghĩa là các chiều đặc trưng hết tương quan với nhau:
\begin{equation}
(DH)(DH)^{\top} = I \;\Longrightarrow\; D^{\top}D = (HH^{\top})^{-1} = G^{-1}.
\end{equation}
Khi đó tích vô hướng giữa đặc trưng đã tẩy trắng và prototype đã tẩy trắng chính
là
\begin{equation}
\hat s_y = (D h_{\text{test}})^{\top}(D c_y) = h_{\text{test}}^{\top} G^{-1} c_y .
\end{equation}

\begin{luuy}
Nói ngắn gọn: \textbf{tẩy trắng dữ liệu rồi mới so sánh}, thay vì so sánh trực
tiếp trên một không gian đã bị méo. Đó là toàn bộ nội dung của $G^{-1}$.
\end{luuy}

\subsection{Ba mối liên hệ lý thuyết đáng nhớ}

\paragraph{Liên hệ với hồi quy ridge.}
Dạng $W_o = (G + \lambda I)^{-1} C$ \textbf{chính xác là nghiệm đóng} của bài
toán bình phương tối thiểu có điều chuẩn $l_2$, khớp nhãn one-hot:
\begin{equation}
W_o = \arg\min_{W}
\left( \|Y_{\text{train}}^{\top} - W^{\top}H\|_2^2 + \lambda\|W\|_2^2 \right).
\end{equation}
Hệ quả thực tiễn rất mạnh: huấn luyện một lớp tuyến tính bằng SGD với mất mát
bình phương và weight decay bằng $\lambda$ thì \emph{giỏi lắm cũng chỉ tiệm cận}
được nghiệm này. Ở đây ta tính thẳng bằng công thức, không cần lặp.

\paragraph{Liên hệ với LDA và khoảng cách Mahalanobis.}
Công thức \eqref{eq:gram} rất giống LDA, nhưng LDA dùng ma trận \emph{hiệp phương
sai} $S$ còn RanPAC dùng ma trận \emph{Gram} $G$. Hai cái chỉ trùng nhau khi mọi
chiều đặc trưng có trung bình bằng $0$, điều nói chung không đúng. Dùng $G$ có
hai cái lợi: không phải tính thành phần bias, và việc cộng dồn qua các tác vụ đơn
giản hơn vì không phải trừ trung bình.

\paragraph{Liên hệ với ZCA whitening.}
Nếu tất cả các lớp có xác suất tiên nghiệm bằng nhau, phân loại LDA tương đương
với việc tẩy trắng ZCA rồi tìm khoảng cách Euclid nhỏ nhất. Công thức
\eqref{eq:gram} là cùng một tinh thần, chỉ khác ở chỗ đưa ma trận Gram về đơn vị
thay vì đưa hiệp phương sai về đơn vị.

% =====================================================================
\section{Ý tưởng thứ hai: chiếu ngẫu nhiên kèm phi tuyến}

Chiếu đặc trưng $L = 768$ chiều lên $M = 10\,000$ chiều bằng ma trận ngẫu nhiên
$W$ đóng băng, rồi cho qua hàm phi tuyến $\phi$ (ReLU):
\begin{equation}
h = \phi(f^{\top} W).
\end{equation}

\subsection{Vì sao phi tuyến là bắt buộc}

Đây là điểm dễ hiểu nhầm nhất, và lập luận rất gọn:

\begin{luuy}
Nếu bỏ $\phi$ đi, phép chiếu ngẫu nhiên chỉ còn là một \textbf{phép đổi cơ sở
tuyến tính}. Mà một bộ phân lớp tuyến tính thì \emph{bất biến} với phép đổi cơ
sở. Chiếu lên một triệu chiều cũng vô ích.
\end{luuy}

Thí nghiệm ở Hình 3(a) của bài xác nhận đúng như vậy: bỏ kích hoạt phi tuyến thì
dù $M$ lớn đến đâu, kết quả cũng không hơn gì không chiếu.

\subsection{Vậy phi tuyến mang lại gì}

Nó tạo ra \textbf{số hạng tương tác giữa các đặc trưng}. Khai triển Taylor của
$\phi$ sinh ra các tích chéo giữa những chiều đã chiếu -- thứ mà một mô hình
tuyến tính trên đặc trưng gốc không bao giờ chạm tới được.

Bài chứng minh điều này bằng một thí nghiệm rất sạch (Hình 3(b)): lấy 100 đặc
trưng đầu tiên, dựng \emph{tường minh} toàn bộ tích chéo từng cặp, rồi phân loại.
Kết quả tốt hơn hẳn đặc trưng thô. Nhưng với $L = 768$ thì số tích chéo tường
minh lên tới gần $300\,000$, bất khả thi về tính toán. Chiếu ngẫu nhiên kèm phi
tuyến là \textbf{cách rẻ để lấy gần hết lợi ích đó}.

\subsection{Vì sao phải mở rộng chiều chứ không nén}

Vì $W$ là ngẫu nhiên chứ không được học, nên cần dư thừa để bù lại. Bài thử nén
từ 768 xuống 500 chiều và kết quả \emph{thấp hơn cả} phương án không dùng chiếu.

\begin{table}[htbp]
\centering
\caption{Độ chính xác trên Split-CIFAR100 theo số chiều chiếu $M$, không dùng
Phase 1, $\lambda$ cố định bằng 100 (Bảng A5 của bài báo).}
\label{tab:scale}
\begin{tabular}{lrrrrrrr}
\toprule
$M$ & 100 & 400 & 1250 & 2500 & 5000 & 10000 & 15000 \\
\midrule
Độ chính xác (\%) & 71.6 & 83.9 & 86.8 & 87.7 & 88.4 & 88.8 & 89.0 \\
\bottomrule
\end{tabular}
\end{table}

Mức tăng bão hòa dần sau khoảng $M = 5000$.

% =====================================================================
\section{Vì sao cách này hợp với học liên tục đến mức hoàn hảo}
\label{sec:coloi}

Đây là điểm cốt lõi của cả bài, và nó rất thanh lịch.

\begin{equation}
G = \sum_{t=1}^{T}\sum_{n=1}^{N_t} h_{t,n} \otimes h_{t,n},
\qquad
C = \sum_{t=1}^{T}\sum_{n=1}^{N_t} h_{t,n} \otimes y_{t,n}.
\label{eq:congdon}
\end{equation}

Cả hai đều là \textbf{tổng cộng dồn}. Mà phép cộng có tính giao hoán, nên:

\begin{luuy}
Học tuần tự qua $T$ tác vụ cho ra kết quả \textbf{giống hệt} học tất cả cùng một
lúc. Không phải ``quên ít'', mà là \textbf{không quên, chứng minh được bằng đại
số}, không cần thí nghiệm. Thứ tự các tác vụ hoàn toàn không ảnh hưởng.
\end{luuy}

Lý do sâu xa: không có bước cập nhật trọng số lặp nào để mà ghi đè kiến thức cũ.
Toàn bộ ``huấn luyện'' của Phase 2 chỉ là cộng ma trận.

Đó cũng là lý do bài nhấn mạnh cụm \emph{training-free}: khả năng chống quên
không đến từ một cơ chế chống quên nào cả, mà đến từ việc \textbf{chọn một dạng
bài toán vốn không có chỗ cho việc quên xảy ra}.

% =====================================================================
\section{Phase 1: vá khoảng cách miền}

Vẫn còn một vấn đề. Đặc trưng đóng băng có thể không hợp với miền dữ liệu mới --
ảnh phong cách hóa, ảnh y tế, ảnh vệ tinh.

Phase 1 xử lý bằng cách huấn luyện một bộ \textbf{PETL} (Parameter-Efficient
Transfer Learning: AdaptFormer, SSF, hoặc VPT, mỗi bộ vài trăm nghìn tham số)
\textbf{chỉ trên tác vụ đầu tiên} $D_1$, rồi đóng băng vĩnh viễn. Đầu phân lớp
tạm dùng để huấn luyện nó bị vứt bỏ trước khi vào Phase 2.

Logic của lựa chọn này gồm hai vế. Thứ nhất, chỉ huấn luyện ở một tác vụ duy nhất
thì không có cơ hội để quên tích lũy. Thứ hai, dữ liệu tác vụ 1 dù ít vẫn đại
diện cho miền mới tốt hơn nhiều so với ImageNet-21K.

Chi tiết huấn luyện: SGD, batch 48, tốc độ học 0.01, weight decay 0.0005,
momentum 0.9, 20 epoch, lịch cosine annealing.

% =====================================================================
\section{Thuật toán đầy đủ}

\begin{figure}[htbp]
\centering
\resizebox{\textwidth}{!}{%
\begin{tikzpicture}[
  font=\footnotesize, >={Stealth[length=2mm]},
  box/.style={draw, minimum height=8mm, minimum width=2.0cm, inner xsep=4pt, align=center},
  frz/.style={box, fill=black!8},
  lbl/.style={font=\scriptsize, align=center}
]
% Phase 1
\node[lbl, anchor=west] at (-0.6, 1.9) {\textbf{Phase 1} — chỉ chạy trên tác vụ đầu tiên};
\node[box] (d1)  at (0.9, 1.0)  {Dữ liệu $D_1$};
\node[box] (petl) at (4.2, 1.0) {Backbone đóng băng\\ + tham số PETL};
\node[box] (tmp) at (7.8, 1.0)  {Đầu phân lớp tạm\\ (vứt bỏ sau đó)};
\draw[->] (d1) -- (petl);
\draw[->] (petl) -- (tmp);

% Phase 2
\node[lbl, anchor=west] at (-0.6, -0.5) {\textbf{Phase 2} — chạy trên toàn bộ chuỗi $D_1 \ldots D_T$};
\node[box] (dt)  at (0.9, -1.6)  {Dữ liệu $D_t$};
\node[frz] (bb)  at (4.2, -1.6)  {Backbone đã thích nghi\\ (đóng băng)};
\node[frz] (rp)  at (7.8, -1.6)  {Chiếu ngẫu nhiên $W$\\ (đóng băng)};
\node[box] (act) at (10.9, -1.6) {ReLU};
\node[box] (acc) at (14.0, -1.6) {Cộng dồn\\ $G$ và $C$};
\node[box] (sol) at (14.0, -3.4) {$W_o = (G+\lambda I)^{-1}C$};
\node[box] (out) at (10.4, -3.4) {Dự đoán lớp};

\draw[->] (dt) -- (bb);
\draw[->] (bb) -- node[lbl, above] {$f$, 768 chiều} (rp);
\draw[->] (rp) -- (act);
\draw[->] (act) -- node[lbl, above] {$h$, 10000 chiều} (acc);
\draw[->] (acc) -- node[lbl, right] {cuối mỗi tác vụ} (sol);
\draw[->] (sol) -- (out);

\draw[->, dashed] (petl.south) -- node[lbl, right] {tham số PETL, đóng băng} (bb.north);
\end{tikzpicture}}
\caption{Hai pha của RanPAC. Các khối tô xám là thành phần đóng băng, không có
tham số nào được huấn luyện. Toàn bộ ``huấn luyện'' của Phase 2 chỉ là cộng dồn
hai ma trận rồi giải một hệ tuyến tính.}
\label{fig:ranpac}
\end{figure}

\paragraph{Phase 1 -- thích nghi phiên đầu (tùy chọn).}
\begin{enumerate}[leftmargin=*]
\item Với mỗi ảnh trong $D_1$, trích đặc trưng và huấn luyện tham số PETL bằng
      SGD với mất mát cross-entropy trên $K_1$ lớp của tác vụ đầu.
\item Đóng băng tham số PETL. Vứt bỏ đầu phân lớp tạm.
\end{enumerate}

\paragraph{Phase 2 -- học liên tục.}
\begin{enumerate}[leftmargin=*]
\item Sinh ma trận chiếu $W \sim \mathcal{N}(0,1)$ kích thước $L \times M$, đóng
      băng vĩnh viễn.
\item Với mỗi tác vụ $t = 1,\ldots,T$ và mỗi ảnh trong $D_t$: trích đặc trưng
      $f$, tính $h = \phi(f^{\top}W)$, rồi cộng dồn $G$ và $C$ theo
      \eqref{eq:congdon}.
\item Cuối mỗi tác vụ: chọn $\lambda$, rồi tính $W_o = (G + \lambda I)^{-1}C$.
\end{enumerate}

\subsection{Cách chọn \texorpdfstring{$\lambda$}{lambda}}

Bước này đáng chú ý vì nó tuân thủ đúng ràng buộc của học liên tục. Ở mỗi tác vụ:

\begin{enumerate}[leftmargin=*]
\item Chia dữ liệu \emph{của chính tác vụ đó} theo tỉ lệ 80:20.
\item Quét $\lambda$ qua 17 bậc độ lớn, từ $10^{-8}$ đến $10^{8}$.
\item Với mỗi $\lambda$, tính $W_o$ từ phần 80\,\% rồi đo sai số bình phương
      trung bình trên phần 20\,\%.
\item Chọn $\lambda$ cho sai số nhỏ nhất.
\end{enumerate}

Không đụng tới dữ liệu của các tác vụ cũ ở bất kỳ bước nào. Bài cũng thẳng thắn
ghi nhận rằng chọn $\lambda$ chỉ dựa trên tác vụ hiện tại có thể chưa tối ưu so
với khi được xem toàn bộ dữ liệu.

% =====================================================================
\section{Kết quả}

Chỉ số báo cáo là $A_T = \frac{1}{T}\sum_{i=1}^{T} R_{T,i}$ -- trung bình độ
chính xác trên mọi tác vụ, đo ở giai đoạn cuối cùng.

\begin{table}[htbp]
\centering
\caption{Kết quả và khảo sát loại bỏ thành phần, $T = 10$ (Bảng 1 và Bảng 2 của
bài báo). Hai hàng in nghiêng là mốc tham chiếu, không phải phương pháp học liên
tục.}
\label{tab:ketqua}
\begin{tabular}{lrrrr}
\toprule
Cấu hình & CIFAR100 & ImageNet-R & ImageNet-A & CUB \\
\midrule
NCM đơn thuần (cosine)              & 83.4 & 61.2 & 49.3 & 85.1 \\
NCM kèm Phase 1                     & 87.8 & 70.1 & 49.7 & 85.4 \\
Gram, không chiếu, không Phase 1    & 87.0 & 67.6 & 54.4 & 86.8 \\
Gram kèm chiếu, không Phase 1       & 89.0 & 71.8 & 58.2 & 88.7 \\
Gram kèm Phase 1, không chiếu       & 90.6 & 73.9 & 57.7 & 86.6 \\
\midrule
\textbf{RanPAC đầy đủ}              & \textbf{92.2} & \textbf{77.9} & \textbf{62.4} & \textbf{90.3} \\
\midrule
\emph{Linear probe huấn luyện chung}      & \emph{87.9} & \emph{72.0} & \emph{56.6} & \emph{88.7} \\
\emph{Fine-tune toàn bộ, huấn luyện chung} & \emph{93.8} & \emph{86.6} & \emph{70.8} & \emph{90.5} \\
\bottomrule
\end{tabular}
\end{table}

Ba điều đọc ra được:

\begin{enumerate}[leftmargin=*]
\item \textbf{Chiếu ngẫu nhiên và Phase 1 bổ trợ nhau, không thay thế nhau.} Bỏ
      chiếu mất 1.6--4 điểm; bỏ Phase 1 mất 2.2--6.1 điểm; bỏ cả hai mất
      5.2--10.3 điểm.
\item \textbf{RanPAC vượt cả linear probe được huấn luyện chung} ở cả bảy bộ dữ
      liệu trong bài. Đây là điều đáng ngạc nhiên: một phương pháp học tuần tự
      lại hơn một mô hình được xem toàn bộ dữ liệu cùng lúc. Lý do là linear probe
      làm việc trên 768 chiều, còn RanPAC làm việc trên 10\,000 chiều có tương
      tác phi tuyến.
\item \textbf{Khoảng cách tới cận trên rất hẹp trên một số tập.} Trên CIFAR100 nó
      chỉ kém fine-tune toàn bộ 1.6 điểm, trên CUB kém 0.2 điểm. Khoảng cách lớn
      chỉ còn ở hai biến thể ImageNet.
\end{enumerate}

Bài cũng chạy trên ba bộ domain-incremental (CORe50, CDDB-Hard, DomainNet) với
cùng xu hướng, và kiểm chứng thêm với backbone ResNet và CLIP.

% =====================================================================
\section{Cái giá phải trả}

\begin{itemize}[leftmargin=*]
\item \textbf{Bộ nhớ khi huấn luyện.} Ma trận Gram kích thước $M \times M$, tức
      $10\,000^2$ phần tử. Đây là khoản lớn nhất. Bài lập luận rằng vẫn nhỏ hơn
      SLCA, vốn lưu một ma trận hiệp phương sai cho \emph{mỗi} lớp.
\item \textbf{Số tham số.} Thêm $L \times M = 7.68$ triệu trọng số đóng băng (có
      thể lưu 1 bit mỗi phần tử nếu dùng ma trận nhị cực), cùng khoảng 2--2.5
      triệu tham số học được với $K = 200$. So sánh: L2P và DualPrompt dùng
      khoảng 0.3--0.5 triệu.
\item \textbf{Tốc độ suy luận.} Gần như không đổi so với backbone gốc, vì lớp
      chiếu và đầu phân lớp đều chỉ là phép nhân ma trận.
\item \textbf{Tốc độ huấn luyện.} Phase 2 chỉ chậm hơn một lượt chạy suy luận
      chút ít. Phần chậm nhất là nghịch đảo ma trận Gram khi quét $\lambda$,
      khoảng 1 phút mỗi tác vụ trên CPU với $M = 10\,000$.
\item \textbf{Phần cứng.} Toàn bộ thí nghiệm của bài chạy trên một GPU
      RTX 4090 duy nhất.
\end{itemize}

% =====================================================================
\section{Hạn chế mà chính tác giả thừa nhận}

Toàn bộ giá trị của phương pháp \textbf{phụ thuộc hoàn toàn vào chất lượng của bộ
trích đặc trưng có sẵn}. Nó không học biểu diễn, nó chỉ khai thác biểu diễn. Vì
vậy nó khó có thể mạnh tương đương trong các thiết lập phải huấn luyện mạng từ
đầu.

Ngoài ra nó dùng nhiều tham số hơn các phương pháp prompting, dù tác giả cho rằng
đánh đổi này xứng đáng với sự đơn giản khi cài đặt.

% =====================================================================
\section{Tóm tắt và liên hệ với đề tài của chúng ta}

\begin{luuy}
\textbf{Một câu tóm tắt cả bài:} RanPAC là \emph{hồi quy ridge chính xác, trên
đặc trưng đóng băng đã được chiếu ngẫu nhiên lên chiều cao kèm phi tuyến}. Và vì
nghiệm ridge chỉ phụ thuộc vào hai tổng cộng dồn, nó bất biến với thứ tự dữ liệu
-- nên bài toán quên thảm khốc không bị \emph{giải quyết}, mà bị định nghĩa cho
biến mất.
\end{luuy}

Trong đề tài của chúng ta, thành phần được lấy từ bài này là \textbf{toàn bộ
Phase 2}: lớp chiếu ngẫu nhiên đóng băng kèm ReLU, ma trận Gram, class-prototype
và nghiệm ridge. Chúng ta \textbf{không} cài Phase 1, mà thay bằng bộ LoRA tác vụ
0 mà HRM-PET vốn đã huấn luyện sẵn -- nhờ vậy không phải huấn luyện thêm một tham
số nào.

Hai điểm khác biệt cần nhớ khi đối chiếu số liệu:

\begin{itemize}[leftmargin=*]
\item RanPAC \textbf{không có bước định tuyến}. Nó là phương pháp một đầu phân
      lớp duy nhất trải trên mọi lớp; trong cả bài không có khái niệm suy luận
      danh tính tác vụ. Việc gộp điểm số theo lớp của nó thành điểm theo tác vụ
      để làm bộ định tuyến là cách dùng không có trong bài gốc.
\item Chỉ số $A_T$ của họ \textbf{trùng định nghĩa} với cột Acc@1 trong bảng kết
      quả của chúng ta, nên hai bên so sánh được về mặt chỉ số. Tuy nhiên seed và
      cách chia lớp khác nhau, và chúng ta chưa từng tự chạy RanPAC, nên chưa thể
      coi là so sánh trực tiếp.
\end{itemize}

\begin{thebibliography}{9}
\bibitem{ranpac}
Mark D. McDonnell, Dong Gong, Amin Parvaneh, Ehsan Abbasnejad, Anton van den
Hengel.
\newblock RanPAC: Random Projections and Pre-trained Models for Continual
Learning.
\newblock \emph{Advances in Neural Information Processing Systems 36 (NeurIPS
2023)}. arXiv:2307.02251.
\end{thebibliography}

\end{document}
